% Vietnamese translation for OI Public Library
% Source-File: IOI中国国家候选队论文/2008/Day2/7.任一恒《非完美算法初探》/非完美算法初探——任一恒.ppt
\documentclass[11pt]{article}
\usepackage{oipl-vn}

\newcommand{\Oh}{\mathcal{O}}

\title{Bản trình chiếu: Bước đầu tìm hiểu thuật toán không hoàn hảo}
\author{Ren Yiheng}
\date{2008}

\begin{document}
\maketitle

\origtitle{Slide companion for imperfect algorithms}
\sourcepath{IOI China candidate team papers / 2008 / Day 2 / Ren Yiheng.ppt}

\section{Phạm vi bản dịch}

Tệp PowerPoint đi cùng bài viết đã được dịch từ tài liệu Word. Bản ghi chú này
giữ lại các nhóm kỹ thuật không hoàn hảo nhưng thực dụng: tham lam, ngẫu
nhiên, lấy mẫu và mô phỏng luyện kim.

\section{Không hoàn hảo nhưng có kiểm soát}

Một thuật toán không hoàn hảo có thể không bảo đảm tối ưu hoặc không bảo đảm
đúng tuyệt đối trên mọi dữ liệu. Điều này không có nghĩa là tùy tiện: ta vẫn
cần biết thuật toán có thể sai ở đâu, xác suất hoặc mức độ sai ra sao, và vì
sao nó có ích trong giới hạn thời gian cho trước.

\section{Các kỹ thuật}

Tham lam cho nghiệm nhanh nhưng có thể mắc phản ví dụ. Ngẫu nhiên giúp tránh
cấu hình xấu hoặc tìm kiếm nhiều vùng nghiệm. Lấy mẫu hữu ích khi tập tốt rất
lớn hoặc không tồn tại. Mô phỏng luyện kim và tìm kiếm cục bộ dùng hàm đánh
giá để cải thiện nghiệm từng bước.

\section{Kết luận}

Trong bài chấm theo điểm hoặc bài quá khó để giải trọn vẹn, thuật toán không
hoàn hảo là một công cụ thực tế. Tuy nhiên, khi bài yêu cầu đúng tuyệt đối,
chỉ nên dùng chúng như thành phần phụ hoặc phương án lấy điểm.

\end{document}
